\chapter{Analysis}
\label{ch:analysis}

This chapter analyzes the two systems along the dimensions that determine whether learned path control is
practical, replacing the security-centric analysis of the original template (as motivated in
\cref{sec:problem:trust}) with the concerns that actually govern deployment: overhead
(\cref{sec:analysis:overhead}), scalability (\cref{sec:analysis:scalability}), generalization
(\cref{sec:analysis:generalization}), and deployability (\cref{sec:analysis:deployability}). Where a claim
depends on measurements still in progress, it is marked.

\section{Overhead}
\label{sec:analysis:overhead}

Two kinds of overhead matter: measurement and computation.

\paragraph{Measurement (probing) overhead.} In Part~I, probing is a first-class cost charged in the reward
(\cref{eq:p1reward}) and accounted per decision (\cref{sec:p1impl:baselines}). The learned selector probes
only the paths it inspects, in contrast to heuristics that require a metric on every path, so its
measurement footprint is intended to be a fraction of an exhaustive scan at comparable selection quality --
the property we test in \cref{sec:p1eval:overhead}. In Part~II, per-path state is derived from ordinary
transport feedback (RTT samples, acknowledgements, loss signals) rather than dedicated probes, so the
marginal measurement cost of the split decision is low.

\paragraph{Computational overhead.} Both policies are small neural networks evaluated once per decision.
The Part~I selector runs at path-selection granularity (seconds to minutes) and is comfortably cheap. The
Part~II transport agent is the tighter constraint: it runs \emph{per frame} (${\approx}\SI{33}{\milli\second}$
at 30\,fps), so its inference must fit well inside a frame interval. The scoring architecture keeps this
affordable -- a shared per-path encoder evaluated $N$ times and a softmax -- but the per-frame budget is the
binding real-time constraint. \gap{report measured per-decision inference latency for both agents, and
confirm the transport agent's inference fits within the frame budget on the target hardware.}

\section{Scalability}
\label{sec:analysis:scalability}

The scoring architecture is what makes both systems scale in the number of paths. Because a shared network
scores each path independently and decisions are made by comparison (\cref{sec:p1design:scoring}) or softmax
(\cref{sec:p2design:scoring}), the parameter count is independent of $N$, and inference is linear in $N$.
This lets one trained model serve source--destination pairs with any number of paths and topologies of any
size, without retraining or padding to a fixed maximum -- the limitation of the flat agents. The remaining
scalability question is training coverage: the environment must expose enough variety of path counts,
topologies, and conditions for the shared scorer to generalize. \gap{report how selection quality holds as
topology size and path count grow, ideally on topologies larger than those seen in training.}

\section{Generalization}
\label{sec:analysis:generalization}

Generalization is required along three axes. \textbf{Across intents:} Part~I's FiLM conditioning is meant to
interpolate to intents between (and, ideally, beyond) the trained reward profiles, so that a novel weight
vector yields sensible behavior without retraining -- tested in \cref{sec:p1eval:intent}. \textbf{Across
topologies and path counts:} the permutation-equivariant design gives structural generalization, but only
empirical results can confirm it holds off the training distribution. \textbf{Across conditions:} the
dynamic scenario of Part~II (\cref{sec:p2eval:results}) is itself a generalization stress test, since the
policy must cope with churn and bursts rather than a stationary environment; the collapse of the
fixed-schedule baselines there shows that static policies do \emph{not} generalize across volatility, which
is the gap learned control fills. \gap{quantify off-distribution degradation for both systems -- unseen
intents, larger topologies, and unseen dynamics regimes -- to bound how gracefully performance decays.}

\section{Deployability}
\label{sec:analysis:deployability}

Several gaps separate these prototypes from a deployment.

\begin{itemize}
  \item \textbf{From simulation to reality.} Both systems are trained in simulation. Part~I relies on a
  synthetic traffic and beaconing environment; Part~II on NS-3 and a mock. A sim-to-real transfer -- domain
  randomization, or online fine-tuning against live measurements -- would be needed before deployment, and
  the packet-level NS-3 backend is the intended stepping stone (\cref{sec:p2impl:gaps}).
  \item \textbf{Real transport.} Part~II models subflows with TCP-like congestion control rather than true
  Multipath QUIC~\cite{DeConinck2017}; a real QUIC integration is required to validate the per-frame
  scheduling under an actual multipath transport.
  \item \textbf{Online and safe learning.} A deployed agent would learn or adapt against live traffic, where
  exploration has real cost. Bounding worst-case behavior -- e.g.\ falling back to a safe heuristic when the
  policy is uncertain, or constraining exploration -- is a prerequisite we do not address here.
  \item \textbf{Trust and integrity.} As noted in \cref{sec:problem:trust}, the agents assume honest
  measurements. A deployment must verify measurement integrity and bound the damage a single misreported
  path can cause; Part~I's trust term is a partial hedge, not a solution.
\end{itemize}

On the positive side, two properties aid deployability: the endpoint-driven nature of SCION path selection
means these controllers need no network-internal changes, and the horizon-agnostic application agent
(\cref{sec:p2design:hierarchy}) is designed to run in unbounded real sessions rather than fixed-length
episodes.
